% Canadian AI 2027 - position paper, short-paper track (6 pages incl. references)
% FORMAT: PMLR (jmlr class), matching the template CAIAC required for
% Canadian AI 2026 (proceedings in PMLR). Swap in the official CAIAC 2027
% template when it is posted; the page count should be nearly identical.
% DOUBLE-BLIND: the 2026 edition used double-blind review. Before
% submitting, replace the \author block with \author{\Name{Anonymous Authors}}.
\documentclass[pmlr]{jmlr}
\usepackage[T1]{fontenc}
\usepackage{array}
\jmlrvolume{}
\jmlryear{2027}
\jmlrworkshop{Canadian AI 2027}
\title{Grounding the AI Scrum Master: Evidence Before Autonomy}
\author{\Name{Hani AbuSharkh} \Email{habusharkh@hct.ac.ae}\\
\addr Higher Colleges of Technology, Al Ain, UAE}
\begin{document}
\maketitle

\begin{abstract}
Agent systems that carry the Scrum Master label are multiplying, yet few of them do the job, and none we found has been judged against what actually happened in real sprints. Meanwhile, work that predicts sprint outcomes from Jira histories stops at prediction and never acts. We argue that the two should meet, and that the meeting point gives a bar that can be tested today with public data: an agent in this role should ground what it says in leakage-controlled models, and should be judged by replaying historical sprints without hindsight, against a model-only arm and an agent-only arm. Two documented failure modes bracket the design space. An ungrounded agent says things the data does not support; an over-grounded one repeats its model and adds nothing. We define the replay protocol, an agreement rate, and a cost-weighted measure of what grounding adds, so that the claim can be checked rather than asserted.
\end{abstract}

\begin{keywords}
LLM agents, Scrum, agile software development, sprint analytics, evaluation, grounding
\end{keywords}

\section{Introduction}

Much of what a Scrum Master does never reaches the tracker. \citet{shastri} name six everyday functions of the role: facilitating, mentoring, negotiating, adapting processes, coordinating and protecting the team. Only two leave much of a trace in Jira. The rest happen in conversation.

Language-model agents have nonetheless started to wear the title. Since 2024, published systems have given an agent the Scrum Master role in a code-generation pipeline \citep{agilecoder}, simulated agile teams with role agents \citep{cognisim}, and turned burndown charts into status reports under the name of an AI Scrum Master \citep{alliata}.

Our position is that the label has run ahead of the evidence, and that a specific, affordable test can bring it back into line. The test is sprint replay: rebuild a project's board as it stood at a chosen moment, give the agent only what was knowable then, let it decide whether and what to say, and score that against what the sprint went on to do. No partner organisation or live team is needed. What is needed is a sensing layer the agent can cite, and the discipline to report how often the agent merely repeats it.

We argue three things: that the label mostly arrives without the job (Section~\ref{sec:gap}); that two documented failure modes bound what a credible system must avoid (Section~\ref{sec:fail}); and that hindsight-free replay with explicit measures of agreement and added value is strict enough to separate judgement from repetition and cheap enough to apply now (Sections~\ref{sec:ground} and~\ref{sec:replay}).

\section{The label and the evidence}
\label{sec:gap}

Systems that use Scrum roles fall into three families. The first organises code generation: AgileCoder's Scrum Master agent reviews a backlog, picks sprint tasks and judges deliverability, and is scored on code benchmarks alone \citep{agilecoder}. The role supplies a vocabulary, not a set of responsibilities. The second simulates agile teams; CogniSim models scaled-agile roles with language-model agents but has no Scrum Master agent at all, and its own future work asks for predictive analytics for delays \citep{cognisim}. The third targets the role directly, and here the evidence is real but small. An action-research study put assistants into daily Scrums and refinement for three teams at a large organisation: about a third of early refinement outputs were wrong because the input data was stale, two of seven members of one team found the assistant annoying, and participants asked that it inform rather than warn \citep{cabrero}. Another study prompts models to write status reports from burndown charts and scores the text with similarity metrics, with no agent and no prediction \citep{alliata}.

Reviews agree on the shape of the field. Software management accounts for under one percent of language-model usages in a review of 395 software-engineering papers \citep{hou}, and a review of Scrum-specific work finds no applications in the implementation or review-and-retrospect phases \citep{peras}. Table~\ref{tab:gap} sets the role's accountabilities against the best evidence we found.

\begin{table}[t]
\centering\footnotesize
\caption{Scrum Master accountabilities against the best evidence we found. Rows without a citation reflect absences reported by the reviews discussed in the text.}
\label{tab:gap}
\begin{tabular}{>{\raggedright\arraybackslash}p{0.34\linewidth} l >{\raggedright\arraybackslash}p{0.34\linewidth}}
\hline
Accountability & Coverage & Best evidence\\
\hline
Effort and story-point estimation & Well covered & Real data, many projects \citep{deepse}\\
Delivery-risk prediction at planning & Moderate & Real Jira histories \citep{tse2018,sprint2vec}\\
In-sprint monitoring & Thin & Checkpoint prediction, no agent \citep{tse2018}\\
Daily Scrum facilitation & Thin & One field site, three teams \citep{cabrero}\\
Impediment detection and removal & Thin to empty & Detection only; no removal studies\\
Review and stakeholder reporting & Near empty & Text-similarity scoring \citep{alliata}\\
Retrospectives, coaching, team health & Near empty & Small author-labelled studies\\
End-to-end agent in the role & Label only & Code benchmarks and demos \citep{agilecoder}\\
\hline
\end{tabular}
\end{table}

Three patterns follow. The more relational the task, the thinner the evidence: of the six functions \citet{shastri} identify, only facilitation and coordination have any. Evaluation habits are weak, with small samples, single runs and author-labelled ground truth, and where run-to-run stability has been measured it is poor \citep{yoon}. And the data is lopsided. Public sprint histories are plentiful, with one dataset holding thousands of sprints and millions of change-log entries \citep{tawos}, while public transcripts of stand-ups, reviews and retrospectives are, as far as we can tell, absent. The half of the role a tracker records is therefore the half in which an agent can be held to account today.

\section{Two ways to fail}
\label{sec:fail}

\paragraph{Ungrounded.} Asked certification-style Scrum questions, current models recall well and interpret poorly; in one evaluation at least one prompt failed on roughly a sixth of the items, with errors clustering on interpretive topics \citep{psm}. Generated sprint plans were judged unfit for direct use, with priority accuracy swinging from 44 to 92 percent across runs \citep{yoon}. In mission-critical risk analysis, retrieval-grounded models had the lowest hallucination rates \citep{esem24}. In this role the characteristic failure is a confident wrong statement, and the characteristic fix is grounding.

\paragraph{Over-grounded.} The opposite failure is less discussed and, for research, more dangerous. Given a predictive model as a tool, agents tend to copy it: in one study they reproduced a graph model's output on 97.6 to 99.2 percent of cases, and stronger backbones deferred more, not less \citep{tool}. An agent that repeats its model is the model with added latency. Any claim of an agent in this role should therefore report its agreement with the raw model and include a model-only arm. A system that agrees with its tool on nearly every case has demonstrated the tool, not judgement.

Neither failure is particular to Scrum. Agents are being placed over classical predictive models in clinical triage, credit decisions and operations, and the same two questions arise each time: when should the agent believe the model, and when is acting on it worth the cost? Sprints are an unusually good place to study both, because outcomes are public, time-stamped, and free of hindsight once the board is rebuilt.

\section{What grounding means here}
\label{sec:ground}

Grounding means that the agent's statements about risk trace back to models fitted on real sprint histories, under an evaluation that does not leak the future. That last condition carries most of the weight. Some sprint-outcome studies report accuracies near or above 0.9 using features visible only once the sprint is under way \citep{perez}, or after resampling with no stated split order \citep{obike}, and leakage taxonomies now list software engineering among the affected fields \citep{leakage}. A stricter pipeline produces lower numbers, and the lower numbers are the trustworthy ones. An agent that cites an inflated model inherits the inflation and adds a confident voice to it.

Grounding also limits what the agent may suggest. In defect prediction, guidance derived from explanation methods often fails to fix what it predicts and differs from what developers actually change \citep{saner}, while restricting suggestions to changes with historical precedent brings them much closer to real practice \citep{counteract}. Carried into sprints, an agent should not propose cutting scope merely because the model's expectation falls; it should propose changes that comparable teams actually made.

What is left for the agent is narrower than the label and more interesting than the model: reading the issue text the tabular model never sees, deciding whether a situation is worth raising, choosing the moment, and wording it so that a team can act without feeling audited. Each of these is measurable.

\section{Sprint replay as the bar}
\label{sec:replay}

\paragraph{Construction.} Public datasets record field-level changes with timestamps \citep{tawos}, which is enough to rebuild a board as it stood at any point in a sprint. For sprint $s$ and checkpoint $c \in \{0, 0.25, 0.5, 0.75\}$ of the planned duration, let $B_{s,c}$ be the board rebuilt from events time-stamped no later than that point. Every arm sees $B_{s,c}$ and nothing later.

\paragraph{Arms.} Three at minimum: the risk model alone ($M$), the language-model agent without the model ($L$), and the grounded agent ($G$). Without $M$, no one can tell whether the agent added anything; without $L$, no one can tell whether grounding did. At each replay point, arm $a$ emits a decision $d_a(s,c) \in \{\textit{warn}, \textit{silent}\}$ and, for $L$ and $G$, the message it would send.

\paragraph{Agreement.} Over the set $P$ of replay points,
\[
A = \frac{1}{|P|} \sum_{(s,c) \in P} \mathbf{1}\big[d_G(s,c) = d_M(s,c)\big],
\]
reported as a first-class result, not a footnote.

\paragraph{Added value.} We borrow the cost structure of alarm-based prescriptive process monitoring \citep{alarm}. A warning costs $\kappa$, a failed sprint costs $F$, and a warning raised at checkpoint $c$ averts a fraction $e(c)$ of that cost, with $e$ falling as the sprint proceeds. Let $y_s = 1$ if sprint $s$ failed, $w_a(s) = 1$ if arm $a$ warned at any checkpoint, and $c^{*}_a(s)$ the first checkpoint at which it did. Then
\[
C_a(s) = \kappa\, w_a(s) + F\, y_s \big(1 - e(c^{*}_a(s))\, w_a(s)\big),
\qquad
\Delta = \min\big(\bar{C}_M, \bar{C}_L\big) - \bar{C}_G ,
\]
where $\bar{C}_a$ is the mean of $C_a(s)$ over sprints. A positive $\Delta$ means the grounded agent beats the better of the two baselines. Because $\kappa/F$ and $e$ differ between teams, $\Delta$ is reported as a curve over the cost ratio rather than as one number. Read together, a high $A$ with $\Delta \approx 0$ says the agent is the model, while $\Delta > 0$ with $A$ well below one is evidence of judgement.

\paragraph{Reliability and calibration.} Language-model arms are run repeatedly with their variance reported, and any probability an arm states is checked for calibration.

\paragraph{What replay cannot do.} Replay measures what an agent would have said, not whether saying it would have changed the sprint; the counterfactual is not in the log. The one controlled study we know of on language-model facilitation found that it surfaced more information, mostly by drawing in quieter participants, yet did not improve decision quality \citep{alsobay}. Replay is a filter, not a verdict: a system that fails it should not reach a team, and one that passes has earned a study, not a deployment.

\section{Limits}

Public sprint data is open-source, and enterprise teams may behave differently. Models may have seen public tracker data in pre-training, which is one more reason for the $L$ arm and for replication outside software. The relational half of the role stays out of reach while no public ceremony corpus exists. Finally, autonomy: developers keep planning and management work at low autonomy levels \citep{choudhuri}, and the one field study in this role reports participants asking to be informed rather than warned \citep{cabrero}. An agent designed to inform and prepare is easier to accept, and easier to evaluate honestly, than one designed to decide.

\section{Conclusion}

The AI Scrum Master exists as a label and a few pilots, not yet as a system whose judgement has been measured against what really happened in sprints, while sprint-prediction work sits on the other side of the same gap, sensing without acting. The join can be tested cheaply: replay real sprints, ground the agent in models whose evaluation does not leak the future, and report both how often the agent repeated its model and what it added when it did not.

\begin{thebibliography}{23}
\bibitem[Alliata et al.(2024)]{alliata} G. Alliata, A. Singhal, and R. Bozagiu. The AI scrum master. In \emph{Agile Processes in Software Engineering and Extreme Programming: Workshops (XP 2024)}, LNBIP 524. Springer, 2024.
\bibitem[Alsobay et al.(2025)]{alsobay} M. Alsobay, D. M. Rothschild, J. M. Hofman, and D. G. Goldstein. Bringing everyone to the table: an experimental study of LLM-facilitated group decision making. arXiv:2508.08242, 2025. Preprint.
\bibitem[Cabrero-Daniel et al.(2024)]{cabrero} B. Cabrero-Daniel et al. Exploring human-AI collaboration in agile: customised LLM meeting assistants. In \emph{Agile Processes in Software Engineering and Extreme Programming (XP 2024)}, LNBIP. Springer, 2024.
\bibitem[Choetkiertikul et al.(2018)]{tse2018} M. Choetkiertikul, H. K. Dam, T. Tran, A. Ghose, and J. Grundy. Predicting delivery capability in iterative software development. \emph{IEEE Transactions on Software Engineering}, 44(6):551--573, 2018.
\bibitem[Choetkiertikul et al.(2019)]{deepse} M. Choetkiertikul, H. K. Dam, T. Tran, T. Pham, A. Ghose, and T. Menzies. A deep learning model for estimating story points. \emph{IEEE Transactions on Software Engineering}, 45(7):637--656, 2019.
\bibitem[Choetkiertikul et al.(2025)]{sprint2vec} M. Choetkiertikul, P. Banyongrakkul, C. Ragkhitwetsagul, S. Tuarob, H. K. Dam, and T. Sunetnanta. Sprint2Vec: a deep characterization of sprints in iterative software development. \emph{IEEE Transactions on Software Engineering}, 51(1):220--242, 2025.
\bibitem[Choudhuri et al.(2026)]{choudhuri} R. Choudhuri, C. Bird, C. Badea, M. Gerosa, and A. Sarma. You shall not pass! Where and why developers draw the line on AI autonomy. arXiv:2607.00533, 2026. Preprint.
\bibitem[Cinkusz and Chudziak(2025)]{cognisim} K. Cinkusz and J. A. Chudziak. Agile software management with cognitive multi-agent systems. In \emph{Proceedings of the 17th International Conference on Agents and Artificial Intelligence (ICAART)}, 2025.
\bibitem[Esposito et al.(2024)]{esem24} M. Esposito, F. Palagiano, V. Lenarduzzi, and D. Taibi. Beyond words: on large language models actionability in mission-critical risk analysis. In \emph{Proceedings of the 18th ACM/IEEE International Symposium on Empirical Software Engineering and Measurement (ESEM)}, 2024.
\bibitem[Fahrenkrog-Petersen et al.(2022)]{alarm} S. A. Fahrenkrog-Petersen, N. Tax, I. Teinemaa, M. Dumas, M. de Leoni, F. M. Maggi, and M. Weidlich. Fire now, fire later: alarm-based systems for prescriptive process monitoring. \emph{Knowledge and Information Systems}, 64(2):559--587, 2022.
\bibitem[Hou et al.(2024)]{hou} X. Hou et al. Large language models for software engineering: a systematic literature review. \emph{ACM Transactions on Software Engineering and Methodology}, 33(8), 2024.
\bibitem[Kapoor and Narayanan(2023)]{leakage} S. Kapoor and A. Narayanan. Leakage and the reproducibility crisis in machine-learning-based science. \emph{Patterns}, 4(9):100804, 2023.
\bibitem[Lee and Lee(2025)]{saner} G. Lee and S. U.-J. Lee. Can we trust the actionable guidance from explainable AI techniques in defect prediction? In \emph{Proceedings of the IEEE International Conference on Software Analysis, Evolution and Reengineering (SANER)}, 2025.
\bibitem[Nguyen et al.(2025)]{agilecoder} M. H. Nguyen, T. P. Chau, P. X. Nguyen, and N. D. Q. Bui. AgileCoder: dynamic collaborative agents for software development based on agile methodology. In \emph{Proceedings of the IEEE/ACM International Conference on AI Foundation Models and Software Engineering (FORGE)}, 2025.
\bibitem[Obike et al.(2025)]{obike} P. G. Obike, V. E. Ekong, and O. U. Obot. A novel scoring-based Agile sprint deliverability prediction and prioritization using machine learning. \emph{International Journal of Computer Applications}, 187(49):40--53, 2025.
\bibitem[Oueslati et al.(2024)]{counteract} K. Oueslati, G. Laberge, M. Lamothe, and F. Khomh. Mining action rules for defect reduction planning. In \emph{Proceedings of the ACM International Conference on the Foundations of Software Engineering (FSE)}, 2024.
\bibitem[Peras et al.(2026)]{peras} D. Peras et al. Artificial intelligence in Scrum: a systematic literature review. \emph{Informatica}, 2026.
\bibitem[P\'erez Castillo et al.(2024)]{perez} Y. J. P\'erez Castillo, S. D. Orantes Jim\'enez, and P. O. Letelier Torres. Sprint management in Agile approach: progress and velocity evaluation applying machine learning. \emph{Information}, 15(11):726, 2024.
\bibitem[Perkusich et al.(2026)]{psm} M. Perkusich et al. Evaluating large language models on professional scrum master assessments. arXiv:2607.00049, 2026. Preprint.
\bibitem[Shastri et al.(2021)]{shastri} Y. Shastri, R. Hoda, and R. Amor. Spearheading agile: the role of the scrum master in agile projects. \emph{Empirical Software Engineering}, 26(1):3, 2021.
\bibitem[Tawosi et al.(2022)]{tawos} V. Tawosi, A. Al-Subaihin, R. Moussa, and F. Sarro. A versatile dataset of agile open source software projects. In \emph{Proceedings of the 19th International Conference on Mining Software Repositories (MSR)}, pages 707--711, 2022.
\bibitem[Wang and Vemuri(2026)]{tool} Z. Wang and P. Vemuri. When the tool decides: LLM agents defer blindly to graph neural network tools, and stronger backbones defer more. arXiv:2606.14476, 2026. Preprint.
\bibitem[Yoon et al.(2025)]{yoon} Y. Yoon, K. Iwan, M. Zwart, X. Qin, H. Lee, and M. Spichkova. Scrum sprint planning: LLM-based and algorithmic solutions. arXiv:2512.18966, 2025. Preprint.
\end{thebibliography}

\end{document}
